% !TeX root = ../technical_report.tex
\chapter{Open Knowledge Format (OKF v0.2) and Master Agent Routing}
\label{ch:agentic_governance_and_okf}

% Introduction
A robust multi-agent ecosystem requires strict epistemological grounding to prevent catastrophic context drift and AI hallucination during autonomous development cycles. The standardisation of documentation into machine-actionable knowledge graphs forms the foundational stratum for agentic governance within the TaxonVision-MLOps repository. This chapter delineates the implementation of the Google Open Knowledge Format (OKF) specification (version 0.2), the associated automated validation pipelines, and the hierarchical multi-agent routing architecture defined in the root governance matrices. The synchronisation of these systems ensures that both human developers and artificial intelligence agents operate within deterministically verified and temporally accurate contextual bounds.

\section{Open Knowledge Format (OKF v0.2) Schema and Trust Tiers}
\label{sec:okf_schema_and_trust_tiers}

The architectural necessity for machine-readable knowledge management in AI-assisted software repositories is paramount. Traditional, unstructured documentation paradigms suffer from rapid temporal decay and contextual ambiguity, rendering them inadequate for deterministic consumption by autonomous agents. The adoption of the Open Knowledge Format (OKF v0.2) mitigates these failure modes by enforcing strict structural invariants across all documentation assets, transforming passive text into queryable semantic vertices within a comprehensive epistemological network.

At the core of the OKF specification lies the mandatory YAML frontmatter block, which imposes rigid metadata invariants upon every Markdown document within the \texttt{docs/} and \texttt{.agents/} namespaces. The schema dictates the inclusion of standard operational keys: \texttt{type} (e.g., Concept, Architecture Document, Agent Role), \texttt{title}, \texttt{status} (e.g., active, deprecated, draft), \texttt{version}, \texttt{description}, \texttt{tags}, \texttt{generated} (a boolean flag designating machine-authored content), and \texttt{verified}. An essential component of this semantic metadata is the \texttt{stale\_after} attribute, which mandates an explicit expiration timestamp strictly formatted as an ISO 8601 string (e.g., \texttt{YYYY-MM-DDTHH:MM:SSZ}). This automated temporal validation mechanism serves a critical function in the prevention of architectural drift; as repository states evolve, documentation approaching its \texttt{stale\_after} threshold is systematically identified, marking decaying knowledge as stale and prompting necessary human or agentic review cycles before the information becomes structurally misaligned with the active codebase.

Furthermore, the OKF specification implements a robust trust tier taxonomy, as delineated in Section 5.3 of the standard, which categorises the epistemological certainty of each document. The taxonomy defines three distinct confidence intervals, allowing agents to algorithmically weight the reliability of their contextual inputs:
\begin{itemize}
    \item \textbf{Human-Reviewed:} The highest trust tier, indicating that the documentation has been manually audited and validated by a domain expert, thereby authorising its use as a primary ground-truth source for critical agent operations and code synthesis.
    \item \textbf{Machine-Confirmed:} An intermediate tier representing knowledge that has been generated or empirically verified by automated continuous integration (CI) pipelines, static analysis tools, or deterministic execution logs, yet lacks qualitative human verification. This tier is often applied to automatically generated API references and dependency manifest summaries.
    \item \textbf{Unverified:} The lowest tier, applied to draft documents, unvalidated architectural proposals, or raw heuristic outputs from large language models, signalling to downstream agentic consumers that the content must be treated with high epistemic caution and should not be used as the sole basis for autonomous system mutation.
\end{itemize}

By embedding these structural invariants and trust taxonomies directly into the file metadata, the TaxonVision-MLOps ecosystem guarantees that autonomous agents can algorithmically evaluate the validity, freshness, and safety of their contextual priors prior to initiating complex refactoring or inference optimisation tasks.

\section{Automated OKF Verification Engine}
\label{sec:okf_verification_engine}

To enforce the OKF v0.2 schema invariants and guarantee the integrity of the knowledge graph, a rigorous automated verification engine is deployed via the \texttt{scripts/validate\_okf.py} implementation. This Python executable operates as a non-negotiable continuous integration (CI) pre-commit interception gateway, systematically halting pull requests when frontmatter schemas or structural conventions violate the established OKF specification. The execution of this audit script, triggered via the mandatory \texttt{pixi run -e dev python scripts/validate\_okf.py} command, ensures that the repository's documentation remains continuously consistent, pristine, and formally verified.

The algorithmic structure of the verification engine involves a multi-stage parsing and validation pipeline. Initially, the script traverses the \texttt{docs/} and \texttt{.agents/} directory trees, isolating valid Markdown files and extracting their YAML frontmatter blocks using strict boundary detection. A stringent YAML parser subsequently evaluates the structural integrity of the extracted metadata, passing the resultant dictionary to a series of specific validator functions. The validation engine ensures the presence of all mandatory keys, enforces stringency on type constraints, and rejects files with malformed or missing attributes.

Crucially, the engine performs rigorous ISO 8601 timestamp validation on the \texttt{stale\_after} field. To facilitate adherence and alleviate manual correction overhead, an automated normalisation subroutine can be invoked (via the \texttt{--fix} argument) to transmute partial date strings into strict UTC ISO 8601 timestamps. Concurrently, the script aggregates the freshness status of the entire bundle, generating an audit report that quantifies the total count of fresh, stale, and non-expiring documents, thereby providing a clear metric for the repository's documentation debt.

Beyond frontmatter invariant validation, the verification engine executes a comprehensive broken cross-reference link checking algorithm. The script parses the Markdown body content using sophisticated regular expressions to identify local and relative hyperlink declarations. Each internal link target is then mathematically resolved against the physical file system topology. If a referenced document does not exist, or if the resolved path attempts to escape the authorised knowledge bundle domains (a critical security invariant), the audit is flagged as a failure.

The following code listing illustrates the core logic for the cross-document Markdown link validation, ensuring that all semantic edges within the knowledge graph point to verifiable nodes.

\begin{listing}[htbp]
\begin{minted}{python}
def _validate_concept_links(
    body: str,
    file_path: Path,
    root_path: Path,
    base_paths: list[Path],
) -> list[str]:
    """Validate cross-document Markdown links to enforce topological integrity."""
    errors: list[str] = []
    for link in MARKDOWN_LINK_PATTERN.findall(body):
        if link.startswith(("http://", "https://", "#")):
            continue

        if link.startswith("/"):
            target = (root_path / link.lstrip("/")).resolve()
        else:
            target = (file_path.parent / link).resolve()

        if not target.exists():
            errors.append(f"Broken Markdown Link: target '{link}' does not exist.")
        elif not any(target.is_relative_to(bp) for bp in base_paths):
            errors.append(f"Security Alert: Link path '{link}' escapes authorized domains.")
    return errors
\end{minted}
\caption{Algorithmic validation of intra-repository Markdown cross-references to ensure knowledge graph connectivity.}
\label{lst:validate_okf_links}
\end{listing}

Through its integration into the pre-commit hook architecture, the verification engine provides an impenetrable safeguard against architectural degradation, ensuring that all agent-facing instructions and system documentation strictly conform to the mathematical rigour of the OKF standard prior to merging into the main branch.

\section{Master Agent Router and Role-Based AI Governance}
\label{sec:master_agent_router_and_governance}

\providecolor{TaxonPrimary}{HTML}{1976D2}
\providecolor{TaxonSecondary}{HTML}{2E7D32}
\providecolor{TaxonAccent}{HTML}{E65100}
\providecolor{TaxonDark}{HTML}{212121}
\providecolor{mgray}{HTML}{F5F5F5}
\providecolor{mgrayborder}{HTML}{616161}

\begin{figure}[htbp]
    \centering
    \resizebox{\textwidth}{!}{%
    \begin{tikzpicture}[
        node distance=1.2cm and 2cm,
        font=\footnotesize,
        >=stealth,
        % Styles
        headerbox/.style={
            rectangle,
            rounded corners=4pt,
            draw=TaxonPrimary!80,
            fill=TaxonPrimary!8,
            thick,
            inner sep=8pt,
            align=left,
            drop shadow={opacity=0.15, shadow xshift=2pt, shadow yshift=-2pt},
            minimum width=6cm
        },
        bodybox/.style={
            rectangle,
            rounded corners=4pt,
            draw=TaxonSecondary!80,
            fill=TaxonSecondary!8,
            thick,
            inner sep=8pt,
            align=left,
            drop shadow={opacity=0.15, shadow xshift=2pt, shadow yshift=-2pt},
            minimum width=6cm
        },
        auditbox/.style={
            rectangle,
            rounded corners=4pt,
            draw=TaxonAccent!80,
            fill=TaxonAccent!8,
            thick,
            inner sep=8pt,
            align=center,
            drop shadow={opacity=0.15, shadow xshift=2pt, shadow yshift=-2pt},
            minimum width=6cm
        },
        rolebox/.style={
            rectangle,
            rounded corners=4pt,
            draw=TaxonDark!80,
            fill=mgray,
            thick,
            inner sep=6pt,
            align=left,
            drop shadow={opacity=0.1, shadow xshift=1pt, shadow yshift=-1pt},
            minimum width=5.5cm
        },
        groupbox/.style={
            rectangle,
            rounded corners=6pt,
            draw=mgrayborder,
            fill=mgray,
            dashed,
            inner sep=12pt
        },
        connector/.style={
            ->,
            thick,
            draw=mgrayborder
        }
    ]

    % Left Panel: OKF v0.2 Knowledge Document Anatomy
    \node[headerbox] (yaml) {\textbf{YAML Frontmatter Schema}\\ \scriptsize $\bullet$ \texttt{title} \\ \scriptsize $\bullet$ \texttt{type} (Guideline, Concept, Architecture, Runbook) \\ \scriptsize $\bullet$ \texttt{status} (stable, draft, deprecated) \\ \scriptsize $\bullet$ \texttt{stale\_after} (ISO 8601 UTC) \\ \scriptsize $\bullet$ \texttt{version}, \texttt{tags}};
    \node[bodybox, below=0.5cm of yaml] (markdown) {\textbf{Document Body}\\ \scriptsize $\bullet$ GitHub-Flavored Markdown \\ \scriptsize $\bullet$ Strict Hyphenation (\texttt{-}) \\ \scriptsize $\bullet$ Unordered lists (\texttt{*}) \\ \scriptsize $\bullet$ Zero trailing whitespace};
    \node[auditbox, below=0.8cm of markdown] (audit) {\textbf{Automated Audit Gate}\\ \scriptsize \texttt{scripts/validate\_okf.py} \\ \scriptsize Verifies parsing integrity, freshness compliance, \\ \scriptsize and frontmatter correctness across \texttt{.agents/}};

    \begin{scope}[on background layer]
        \node[groupbox, fit=(yaml)(markdown)(audit), label={[font=\bfseries, text=TaxonPrimary]above:OKF v0.2 Knowledge Document Anatomy}] (okf_group) {};
    \end{scope}

    \draw[connector] (yaml) -- (markdown);
    \draw[connector] (markdown) -- (audit);

    % Right Panel: AI Role Decision Matrix & Path Routing
    \node[rolebox, right=2.5cm of yaml.north east, anchor=north] (r1) {\textbf{1. \texttt{@orchestrator}}\\ \scriptsize PM \& OKF enforcement};
    \node[rolebox, below=0.15cm of r1] (r2) {\textbf{2. \texttt{@taxonomic-architect}}\\ \scriptsize DarwinCore schema \& taxon tree};
    \node[rolebox, below=0.15cm of r2] (r3) {\textbf{3. \texttt{@mlops-data-engineer}}\\ \scriptsize Arrow streaming \& DVC S3 pipelines};
    \node[rolebox, below=0.15cm of r3] (r4) {\textbf{4. \texttt{@model-engineer}}\\ \scriptsize Timm backbone \& loss implementations};
    \node[rolebox, below=0.15cm of r4] (r5) {\textbf{5. \texttt{@inference-optimizer}}\\ \scriptsize ONNX Runtime INT8 \& CAM interpretability};
    \node[rolebox, below=0.15cm of r5] (r6) {\textbf{6. \texttt{@uncertainty-safety-engineer}}\\ \scriptsize Conformal prediction \& Helmholtz energy OOD};

    \node[rolebox, right=0.3cm of r1] (r7) {\textbf{7. \texttt{@active-learning-engineer}}\\ \scriptsize Triage queue \& human review};
    \node[rolebox, below=0.15cm of r7] (r8) {\textbf{8. \texttt{@observability-engineer}}\\ \scriptsize Prometheus metrics \& Wasserstein drift};
    \node[rolebox, below=0.15cm of r8] (r9) {\textbf{9. \texttt{@api-and-ui-developer}}\\ \scriptsize FastAPI REST \& HTMX views};
    \node[rolebox, below=0.15cm of r9] (r10) {\textbf{10. \texttt{@qa-automator}}\\ \scriptsize Testing matrix \& code invariants};
    \node[rolebox, below=0.15cm of r10] (r11) {\textbf{11. \texttt{@docs-librarian}}\\ \scriptsize LaTeX technical report \& OKF graph};
    \node[rolebox, below=0.15cm of r11] (r12) {\textbf{12. \texttt{@git-operator}}\\ \scriptsize GPG signed commits \& releases};

    \begin{scope}[on background layer]
        \node[groupbox, fit=(r1)(r6)(r7)(r12), label={[font=\bfseries, text=TaxonDark]above:AI Role Decision Matrix \& Path Routing (\texttt{.agents/AGENTS.md})}] (role_group) {};
    \end{scope}

    % Connection between OKF and Roles
    \draw[connector, dashed, <->] (okf_group.east) -- (role_group.west) node[midway, above, font=\scriptsize, align=center] {Context \& \\ Enforcement};

    \end{tikzpicture}%
    }
    \caption{Open Knowledge Format (OKF v0.2) architecture and path-scoped AI agent role routing matrix.}
    \label{fig:okf_agentic_governance}
\end{figure}


The orchestration of large language models and autonomous execution agents within the repository codebase is governed by a highly structured, role-based delegation matrix, primarily instantiated in the \texttt{AGENTS.md} root configuration file and the \texttt{.agents/index.md} directory index. This Master Agent Router protocol dictates the behavioural constraints, operational boundaries, and contextual scoping for all autonomous entities interacting with the repository.

The necessity for a deterministic routing mechanism stems from the inherent risks of unbounded token consumption and hallucinatory degradation prevalent in unconstrained AI environments. To mitigate these risks, the ecosystem enforces strict path-scoped roles. Rather than permitting an agent to recursively crawl the entire repository -- a process that incurs significant computational waste and dilutes the operational context -- the routing matrix maps specific file paths and architectural domains to specialised agent personas.

As documented in \texttt{.agents/index.md}, the decision matrix assigns distinct operational roles based upon the target modification area, ensuring maximum alignment between the agent's context and the specific engineering task:
\begin{itemize}
    \item \textbf{\texttt{@taxonomic-architect}:} Assigned to modifications within \texttt{src/taxon\_vision/domain/} or \texttt{config/taxa\_catalog.yaml}, overseeing DarwinCore taxonomic hierarchy implementations and metadata schema validations.
    \item \textbf{\texttt{@mlops-data-engineer}:}\\Responsible for \texttt{src/taxon\_vision/data/} and DVC dataset pipeline operations, ensuring strict adherence to Apache Arrow streaming protocols, cryptographic data hashing, and rigorous licensing audits.
    \item \textbf{\texttt{@model-engineer}:} Governs the \texttt{src/taxon\_vision/models/} directory and training loop mechanics, implementing complex biological long-tail Pareto distributions and Class-Balanced Loss algorithms.
    \item \textbf{\texttt{@inference-optimizer}:} Tasked with \texttt{src/taxon\_vision/inference/} and ONNX export routines, ensuring optimal INT8/FP16 quantisation for runtime acceleration on NVIDIA hardware.
    \item \textbf{\texttt{@uncertainty-safety-engineer}:} Handles the \texttt{src/taxon\_vision/uncertainty/} modules, deploying epistemic safety measures such as Split Conformal Prediction for marginal coverage and energy-based out-of-distribution (OOD) detection.
    \item \textbf{\texttt{@git-operator}:} Manages workflow manifests in \texttt{.github/} and dependency definitions within \texttt{pyproject.toml} or \texttt{pixi.toml}, operating largely without strict pre-commit gating mechanisms.
\end{itemize}

These roles, instantiated as discrete Markdown files within the \texttt{.agents/roles/} directory, constitute the definitive operational context for an agent session. The token guardrails mandate that an agent must first consult the quick-route decision matrix, select the singular appropriate role file, and refrain from blindly pre-loading unrelated role instructions or historical memory logs unless explicitly prompted to access the canon.

Furthermore, the Master Agent Router enforces bifurcated commit protocols depending upon the execution environment. Local workstation agents (such as Antigravity) and human operators are subjected to a strict cryptographic signing requirement; all local commits must be GPG or SSH signed (via the \texttt{git commit -S} flag) to establish a non-repudiable lineage of code authorship. Conversely, cloud-based sandbox agents (such as Jules) operating via automated pull request (PR) workflows are expressly forbidden from attempting local GPG signing. Instead, these agents are instructed to push their work to standard PR feature branches without the \texttt{-S} flag, deferring the cryptographic signature process to the GitHub Web-Flow engine upon the final PR merge. This strategic segregation of commit protocols ensures seamless CI/CD integration whilst maintaining an unbroken chain of cryptographic trust across the entire repository.

\section{Interactive D3 Knowledge Graph Visualisation}
\label{sec:interactive_d3_knowledge_graph}

To provide both human operators and diagnostic algorithms with a holistic, macro-level view of the epistemological network, the TaxonVision-MLOps repository features an interactive knowledge graph visualisation engine, implemented within \texttt{scripts/visualize\_okf.py}. This utility parses the OKF v0.2 frontmatter metadata across all documentation assets and dynamically constructs a node-link data structure suitable for rendering via the Cytoscape.js library and D3 force-directed simulation layouts.

The extraction algorithm operates by recursively traversing the \texttt{docs/} and \texttt{.agents/} directories, systematically extracting the semantic identifiers, document types, and trust tiers from the YAML frontmatter. Concurrently, the Markdown bodies are parsed to extract intra-repository hyperlinks, which constitute the directed edges of the knowledge graph. This harvested data is synthesised into a structured JSON payload consisting of a \texttt{nodes} array and an \texttt{edges} array, carefully filtering out any topological edges that target non-existent or external resources.

The resultant standalone HTML artifact (\texttt{viz.html}) leverages a sophisticated force-directed layout algorithm to spatially distribute the nodes, simulating repulsive physical forces between non-connected entities and attractive spring forces along the defined edges. This topological projection provides immediate visual feedback regarding the structural cohesion of the documentation. Clusters of highly interconnected concepts emerge naturally, whilst orphaned documents or excessive centralisation bottlenecks become visually apparent, signalling potential architectural debt that requires remediation.

The visualisation engine incorporates dynamic styling heuristics derived directly from the OKF schema invariants. The visual encoding specifications define unique colour palettes for distinct document types, allowing for rapid categorisation during human review.

\begin{table}[htbp]
    \centering
    \caption{Knowledge Graph Node Taxonomy and Visual Encoding Specifications.}
    \label{tab:knowledge_graph_styling}
    \begin{tabularx}{\textwidth}{l l X}
        \toprule
        \textbf{Document Type} & \textbf{Primary Fill Colour} & \textbf{Semantic Role} \\
        \midrule
        Attested Computation & \texttt{\#FF5722} (Deep Orange) & Mathematical proofs, loss function derivations, and formal model verifications. \\
        Scientific Model & \texttt{\#00E5FF} (Cyan) & Neural network architectures, vision backbones, and conformal algorithms. \\
        Architecture Document & \texttt{\#B388FF} (Purple) & High-level system design, infrastructure topology, and deployment manifests. \\
        Agent Role & \texttt{\#FFD600} (Yellow) & Behavioural instructions and contextual constraints for AI execution routing. \\
        Behavioural Rule & \texttt{\#FF1744} (Red) & Mandatory operational constraints, formatting standards, and CI/CD invariants. \\
        Directory Index & \texttt{\#78909C} (Blue Grey) & Navigational hubs and repository entry points for hierarchical file routing. \\
        \bottomrule
    \end{tabularx}
\end{table}

This interactive visual interface, integrated seamlessly into the final Zensical documentation build, empowers operators to traverse the semantic connections between complex taxonomic algorithms, Kubernetes orchestration manifests, and multi-agent delegation rules. By clicking upon any node within the force-directed simulation, a detailed side-panel reveals the underlying OKF frontmatter, providing immediate access to the trust tier, freshness status, description, and raw asset path of the target document. The force-directed knowledge graph thereby transforms the static documentation repository into a living, navigable representation of the TaxonVision-MLOps cognitive architecture, facilitating both rapid human comprehension and deterministic agentic navigation.
